\paragraph{Endpoint-singular coefficients and an open-slab criterion.}
\label{par:openslab}
The preceding closed-slab criterion includes zero simplicity
because it requires a finite comparison coefficient \emph{at}
$s=T=1/2$. That condition is unnecessarily strong for
propagating nonnegativity to the terminal time.
Indeed, the maximum-principle argument may be applied
on each shorter interval $[0,T']$, $T'<T$, followed by
continuity at $T$. This distinction matters precisely
when real zeros collide at the endpoint.

\begin{theorem}[Preterminal spatial comparison]\label{thm:openslab}
Let $T>0$, and suppose $U$ is bounded and continuous
on $[0,T]\times\mathbb R$ and belongs to
$C^{1,2}((0,T)\times\mathbb R)$.
Suppose $U(0,x)\ge0$ for all real $x$.
If there is a function
$c:[0,T)\times\mathbb R\to[0,\infty)$,
locally bounded on every
$[0,T']\times[-R,R]$ with $T'<T$,
such that
\begin{equation}\label{eq:openslabineq}
U_s-\tfrac12U_{xx}+c(s,x)U\ge0
\qquad(0<s<T,\ x\in\mathbb R),
\end{equation}
then $U(s,x)\ge0$ for \emph{all}
$0\le s\le T$ and $x\in\mathbb R$.
No bound or finite value of $c$ at $s=T$ is required.
\end{theorem}
\begin{proof}
For each $T'\in(0,T)$, apply the locally
bounded-coefficient maximum principle
(Theorem~\ref{thm:spatialmax}) to $U$
on the closed subslab $[0,T']\times\mathbb R$.
It follows that $U\ge0$ for every $s<T$.
Letting $s\uparrow T$ and using continuity
proves the terminal inequality.
\end{proof}

\begin{corollary}[A zero-simplicity-free sufficient
condition]\label{cor:openslabRH}
For the genuine theta heat-flow hierarchy,
put $T=1/2$ and
$U_n(s,x)=L_{n,\,T-s}(x)$.
Suppose that for every $n\ge1$ there is
a nonnegative function $c_n(s,x)$,
locally bounded on compact subsets of
$[0,T)\times\mathbb R$, satisfying
\begin{equation}\label{eq:openslabrh}
(n+1)(2n+1)U_{n+1}(s,x)
\le c_n(s,x)U_n(s,x)
\quad(0\le s<T,\ x\in\mathbb R).
\end{equation}
Then RH holds. This sufficient condition does
\emph{not} additionally require that all
zeros of $H_0$ be simple.
\end{corollary}
\begin{proof}
The exact heat hierarchy gives
$U_{n,s}-\tfrac12U_{n,xx}+c_nU_n\ge0$.
The initial $U_n(0,x)$ are nonnegative
by the classical real-zero theorem for $H_{1/2}$.
Theorem~\ref{thm:openslab}
gives $L_{n,0}(x)\ge0$ for every $n\ge1$
and every real $x$. The standard
Laguerre--P\'olya real-zero characterization
for the classical $\Xi$ entire function
then implies RH. No simplicity conclusion
about $H_0$ is used.
\end{proof}

\begin{proposition}[The open-slab criterion has
precisely RH strength]\label{prop:openslabequivalence}
Using the classical de Bruijn--Newman zero-dynamics
and the Rodgers--Tao lower bound, the existence
assertion in Corollary~\ref{cor:openslabRH}
is equivalent to RH.
\end{proposition}
\begin{proof}
The forward implication is the corollary.
Conversely suppose RH. Then
$\Lambda\le0$, whereas Rodgers and Tao proved
$\Lambda\ge0$; thus $\Lambda=0$.
At every strictly positive heat time $t>\Lambda$
the classical de Bruijn--Newman theory gives
real, simple zeros of $H_t$.
The real-zero canonical product, with infinitely
many zeros, gives
$L_{n,t}(x)>0$ for each fixed $n\ge1$ and real $x$.
Therefore the exact quotient
\[
c_n(s,x)=
(n+1)(2n+1)\,
\frac{L_{n+1,\,T-s}(x)}{L_{n,\,T-s}(x)}
\]
is nonnegative and continuous on $[0,T)\times
\mathbb R$, hence locally bounded on every
compact preterminal subslab. It satisfies
\eqref{eq:openslabrh} with equality.
We make no assertion that this quotient
is finite as $s\uparrow T$.
\end{proof}

\begin{example}[A harmless terminal collision]\label{ex:heatcollision}
The polynomial $H_t(z)=(z-a)^2-2t$ solves the
same backward heat equation,
$\partial_tH_t=-\partial_z^2H_t$.
Its first Laguerre coefficients are
\[
L_{1,t}(x)=2(x-a)^2+4t,\qquad
L_{2,t}(x)=1,\qquad L_{n,t}=0\ (n\ge3).
\]
For $t>0$ the optimal nonnegative
comparison coefficient is
\begin{equation}\label{eq:collisioncost}
c_1(t,x)=
\frac6{2(x-a)^2+4t}
=\frac3{(x-a)^2+2t}.
\end{equation}
It is bounded on every compact positive-time
subslab, but
$c_1(t,a)=3/(2t)\to+\infty$ as $t\downarrow0$.
At $t=0$ the transform has a double real zero,
and the \emph{closed}-slab pointwise comparison
is impossible at $x=a$.
Nevertheless
$L_{1,t}(x)\ge0$ for all $t\ge0$.
Thus this endpoint blow-up invalidates only
the unnecessarily strong closed-slab criterion:
it does not prevent the open-slab maximum
principle from reaching the terminal time.
\end{example}

\paragraph{The actual missing estimate.}
The open-slab reformulation is mathematically
stronger as a proof \emph{method}, because it
does not require zero simplicity at $t=0$.
It is not an independent proof of RH:
the existence of its coefficients is itself
equivalent to RH, and the quotient used in
the converse already assumes real-rootedness.
To make progress one must derive, \emph{without}
assuming real zeros or positivity of the
Laguerre denominators, a theta-specific
comparison on every compact positive-heat-time
subslab. The ratio's possible growth near
$t=0$ need not be bounded uniformly.
A proof that starts from this quotient
without an independent theta estimate
would be circular.
